\documentclass[a4paper]{llncs}
\usepackage{longtable}
\newcommand{\tocTitle}[2]{\protect\contentsline{title}{#1}{#2}}
\newcommand{\tocAuthors}[1]{{\raggedright \leftskip 15pt \rightskip 2.55em\itshape #1\endgraf}}
\newcommand{\tocSection}[1]{\subsection*{#1}}
\pagestyle{plain}
\pagenumbering{roman}
\begin{document}
\setcounter{page}{1}
\pagestyle{empty}

\noindent
{\Large Camilo Rocha (Ed.) }

\vspace{4cm}

\noindent
{\Huge\bf Rewriting Logic and\\ its Applications}

\vspace{1cm}

\noindent
{\LARGE 16th International Workshop, {\sc WRLA 2026} \\[2ex]
part of the European Joint Conferences on\\ Theory \& Practice of Software (ETAPS 2026)
 \\[2ex]
Turin, Italy, April 11th \& 12th, 2026.}

\vspace{1cm}
\begin{center}
{\LARGE Informal Proceedings}
\end{center}

\vfill


\cleardoublepage

\pagestyle{plain}

\section*{Preface}

This volume contains the preliminary proceedings of the 16th
International Workshop on Rewriting Logic and its Applications (WRLA
2026).
%
Previous editions of the workshop were held in Asilomar (USA) 1996,
Pont-\`a-Mousson (France) 1998, Kanazawa (Japan) 2000, Pisa (Italy)
2002, Barcelona (Spain) 2004, Vienna (Austria) 2006, Budapest
(Hungary) 2008, Paphos (Cyprus) 2010, Tallinn (Estonia) 2012, Grenoble
(France) 2014, Eindhoven (Netherlands) 2016, Thessaloniki (Greece)
2018, Dublin (Ireland) 2020, Munich (Germany) 2022, and Luxembourg
(Luxembourg) 2024.

Rewriting logic is a natural model of computation and an expressive
semantic framework for concurrency, parallelism, communication, and
interaction. It supports the specification of a wide range of systems
and languages across diverse application domains and it also exhibits
desirable properties as a metalogical framework for representing
logics. Over the years, several languages based on rewriting logic
have been designed and implemented. The aim of this workshop is to
bring together researchers with a shared interest in rewriting logic
and its applications, providing a forum to present recent work,
discuss future research directions, and exchange ideas.

WRLA 2026 was held on April 11--12, 2026, in Turin, Italy, as a
satellite event of the European Joint Conferences on Theory {\&}
Practice of Software (ETAPS).
%
A total of 15 submissions were received, covering a wide range of
topics related to rewriting logic and its applications. Each
submission was reviewed by at least three program committee members
or additional reviewers. After extensive discussion, the program
committee selected 11 papers for presentation at the workshop.
%
A selection of the accepted papers will appear in the post-proceedings,
to be published in the Springer LNCS series, following the tradition
of previous editions of the workshop.

Sincere appreciation is extended to all authors who submitted papers
to the workshop, the invited speakers, and those who contributed
tutorials. Equal thanks are due to the members of the program
committee and the external reviewers for their careful work throughout
the review process.
%
As in previous editions, the suggestions provided by the members of
the WRLA steering committee were both valuable and greatly
appreciated.
%
Finally, deep gratitude is expressed to the local organizers of ETAPS
2026, whose efforts made this workshop possible.


~\bigskip


\noindent
\begin{minipage}[t]{.4\textwidth}
April, 2026
\end{minipage}%
\hfill
\begin{minipage}[t]{.4\textwidth}\flushright
Camilo Rocha
\end{minipage}

\clearpage


%% \section*{Table of Contents}
%%   \tocTitle{Combining Parallel Graph Rewriting and Quotient Graphs}{1}
%%   \tocAuthors{Thierry Boy de La Tour and Rachid Echahed}
%%   \tocTitle{Connecting Constrained Constructor Patterns and Matching Logic}{16}
%%   \tocAuthors{Xiaohong Chen, Dorel Lucanu and Grigore Roșu}
%%   \tocTitle{Analysis of the Runtime Resource Provisioning of BPMN Processes using Maude}{33}
%%   \tocAuthors{Francisco Dur\'an, Camilo Rocha and Gwen Sala\"un}
%%   \tocTitle{Evaluation of Logic Programs with Built-Ins and Aggregation: A Calculus for Bag Relations}{49}
%%   \tocAuthors{Matthew Francis-Landau, Tim Vieira and Jason Eisner}
%%   \tocTitle{Well-Founded Induction via Term Refinement in CafeOBJ}{64}
%%   \tocAuthors{Kokichi Futatsugi}
%%   \tocTitle{A Rule-based System for Computation and Deduction in Mathematica}{79}
%%   \tocAuthors{Mircea Marin, Besik Dundua and Temur Kutsia}
%%   \tocTitle{Compositional Verification in Rewriting Logic}{94}
%%   \tocAuthors{\'Oscar Mart\'{\i}n, Alberto Verdejo and Narciso Mart\'{\i}-Oliet}
%%   \tocTitle{Variants in the Infinitary Unification Wonderland}{109}
%%   \tocAuthors{Jos\'e Meseguer}
%%   \tocTitle{Variant Satisfiability of Parameterized Strings}{124}
%%   \tocAuthors{Jos\'e Meseguer}
%%   \tocTitle{Inductive Reasoning with Equality Predicates, Contextual Rewriting and  Variant-Based Simplification}{139}
%%   \tocAuthors{Jos\'e Meseguer and Stephen Skeirik}
%%   \tocTitle{A Remark for the Use of a Path Ordering with an Algebra and a Howard-Style Interpretation of Lambda for Termination Proofs of Typed Rewrite Systems}{157}
%%   \tocAuthors{Mitsuhiro Okada and Yuta Takahashi}
%%   \tocTitle{Strategies, model checking and branching-time properties in Maude}{172}
%%   \tocAuthors{Rub\'en Rubio, Narciso Mart\'{\i}-Oliet, Isabel Pita and Alberto Verdejo}
%%   \tocTitle{Verification of the IBOS Browser Security Properties in Reachability Logic}{187}
%%   \tocAuthors{Stephen Skeirik, Jos\'e Meseguer and Camilo Rocha}
%%   \tocTitle{Automated Construction of Security Wrappers for Industry 4.0 Applications}{205}
%%   \tocAuthors{Carolyn Talcott and Vivek Nigam}
%%   \tocTitle{Maude-SE: a Tight Integration of Maude and SMT Solvers}{220}
%%   \tocAuthors{Geunyeol Yu and Kyungmin Bae}

\clearpage


\section*{Program Committee}
\noindent
\begin{longtable}{p{0.35\textwidth}p{0.65\textwidth}}
Kyungmin Bae & Pohang University of Science and Technology (POSTECH), South Korea
\\Francisco Durán & Universidad de Málaga, Spain
\\Santiago Escobar & Universitat Politècnica de València, Spain
\\Nao Hirokawa & JAIST, Japan
\\Alexander Knapp & Universität Augsburg, Germany
\\Temur Kutsia & RISC, Johannes Kepler University Linz, Austria
\\Alberto Lluch-Lafuente & Technical University of Denmark, Denmark
\\Dorel Lucanu & Alexandru Ioan Cuza University, Romania
\\Salvador Lucas & Universitat Politècnica de València, Spain
\\Narciso Martí-Oliet & Universidad Complutense de Madrid, Spain
\\José Meseguer & University of Illinois at Urbana-Champaign, USA
\\César Muñoz & NASA Langley, USA
\\Kazuhiro Ogata & JAIST, Japan
\\Carlos Olarte & Université Sorbonne Paris Nord, France
\\Peter Ölveczky & University of Oslo, Norway
\\Carlos Ramírez & Pontificia Universidad Javeriana, Colombia
\\Adrián Riesco & Universidad Complutense de Madrid, Spain
\\Christophe Ringeissen &  INRIA, France
\\Camilo Rocha & Pontificia Universidad Javeriana, Colombia
\\Rubén Rubio & Universidad Complutense de Madrid, Spain
\\Carolyn Talcott & SRI International, USA
\\
\end{longtable}

\clearpage


\section*{Additional Reviewers}
\begin{longtable}{l}
\\
\textbf{A} \\*
Arusoaie, Andreea-Valentina 
\\
\\
\textbf{B} \\*
Bonchi, Filippo
\\
\\
\textbf{F} \\*
Frey, Jonas
\\
\\
\textbf{S} \\*
Sobocinski, Pawel
\end{longtable}

\end{document}
